% FAC-EN-U0043
% French source: fac_body.tex lines 1549--1571.
% Authority: PDF physical pp. 42--43; printed pp. 237--238 (no. 45 complete).

\medskip
\phantomsection\label{fac:II:3:no:45}
\textbf{45. Sections of a coherent algebraic sheaf on an affine variety.}

\phantomsection\label{fac:II:3:no:45:thm:2}
\textsc{Theorem 2.} \textit{Let $\mathcal{F}$ be a coherent algebraic sheaf on an affine variety $X$. For every $x\in X$, the $\mathcal{O}_{x,X}$-module $\mathcal{F}_x$ is generated by the elements of $\Gamma(X,\mathcal{F})$.}

Since $X$ is affine, we may embed it as a closed subvariety of an affine space $K^r$. Extending the sheaf $\mathcal{F}$ by zero outside $X$, we obtain a coherent algebraic sheaf on $K^r$ (cf.~\hyperref[fac:II:2:no:39]{no.~39}), and are reduced to proving the theorem for this new sheaf. In other words, we may suppose that $X=K^r$.

By the definition of coherent sheaves, there is a covering of $X$ by open sets over which $\mathcal{F}$ is isomorphic to a quotient of a sheaf $\mathcal{O}^p$. Using \hyperref[fac:II:3:no:42:prop:3]{Proposition~3}, we see that there are finitely many polynomials $Q_i$ having no common zero and, over each $U_i=X_{Q_i}$, a surjective homomorphism
\[
\varphi_i:\mathcal{O}^{p_i}\longrightarrow\mathcal{F}.
\]
We may moreover suppose that none of these polynomials is identically zero.

The point $x$ belongs to one of the $U_i$, say $U_0$. It is clear that $\mathcal{F}_x$ is generated by the sections of $\mathcal{F}$ over $U_0$. Since $Q_0$ is invertible in $\mathcal{O}_x$, it is therefore enough to prove the \hyperref[fac:II:3:no:45:lemma:1]{following lemma}.

\phantomsection\label{fac:II:3:no:45:lemma:1}
\textsc{Lemma 1.} \textit{If $s_0$ is a section of $\mathcal{F}$ over $U_0$, there are an integer $N$ and a section $s$ of $\mathcal{F}$ over $X$ such that}
\[
s=Q_0^N s_0
\qquad\textit{over $U_0$.}
\]

By \hyperref[fac:II:3:no:42:prop:2]{Proposition~2}, $U_i\cap U_0$ is an affine variety, and is clearly irreducible. Applying \hyperref[fac:II:3:no:44:prop:7:cor:2]{Corollary~2 of Proposition~7} to this variety and to
\[
\varphi_i:\mathcal{O}^{p_i}\longrightarrow\mathcal{F},
\]
we find a section $\sigma_{0i}$ of $\mathcal{O}^{p_i}$ over $U_i\cap U_0$ such that
\[
\varphi_i(\sigma_{0i})=s_0
\qquad\text{over }U_i\cap U_0.
\]
Now $U_i\cap U_0$ is the locus of the points of $U_i$ at which $Q_0$ does not vanish. We may therefore apply \hyperref[fac:II:3:no:43:prop:5]{Proposition~5} with $X=U_i$ and $Q=Q_0$. Thus, for every sufficiently large $n$, there is a section $\sigma_i$ of $\mathcal{O}^{p_i}$ over $U_i$ which agrees with $Q_0^n\sigma_{0i}$ over $U_i\cap U_0$. Put
\[
s_i'=\varphi_i(\sigma_i).
\]
Then $s_i'$ is a section of $\mathcal{F}$ over $U_i$ which agrees with $Q_0^n s_0$ over $U_i\cap U_0$.

The sections $s_i'$ and $s_j'$ agree on $U_i\cap U_j\cap U_0$. Applying \hyperref[fac:II:3:no:43:prop:6]{Proposition~6} to $s_i'-s_j'$, we see that, for every sufficiently large $m$,
\[
Q_0^m(s_i'-s_j')=0
\qquad\text{on all of }U_i\cap U_j.
\]
The sections $Q_0^m s_i'$ consequently define a unique section $s$ of $\mathcal{F}$ over $X$, and
\[
s=Q_0^{n+m}s_0
\qquad\text{over }U_0.
\]
This proves the lemma and completes the proof of \hyperref[fac:II:3:no:45:thm:2]{Theorem~2}.

\phantomsection\label{fac:II:3:no:45:thm:2:cor:1}
\textsc{Corollary 1.} \textit{The sheaf $\mathcal{F}$ is isomorphic to a quotient sheaf of some $\mathcal{O}_X^p$.}

Since $\mathcal{F}_x$ is a finitely generated $\mathcal{O}_{x,X}$-module, the \hyperref[fac:II:3:no:45:thm:2]{theorem} shows that there are finitely many sections of $\mathcal{F}$ which generate $\mathcal{F}_x$. By \hyperref[fac:I:2:no:12:prop:1]{Proposition~1 of no.~12}, these sections also generate $\mathcal{F}_y$ when $y$ is sufficiently near $x$. Since $X$ is quasi-compact, we conclude that there are finitely many sections
\[
s_1,\ldots,s_p
\]
of $\mathcal{F}$ which generate $\mathcal{F}_x$ for every $x\in X$. This says precisely that $\mathcal{F}$ is isomorphic to a quotient sheaf of $\mathcal{O}_X^p$.

\phantomsection\label{fac:II:3:no:45:thm:2:cor:2}
\textsc{Corollary 2.} \textit{Let}
\[
\mathcal{A}\xrightarrow{\ \alpha\ }\mathcal{B}
\xrightarrow{\ \beta\ }\mathcal{C}
\]
\textit{be an exact sequence of coherent algebraic sheaves on an affine variety $X$. Then the sequence}
\[
\Gamma(X,\mathcal{A})\xrightarrow{\ \alpha\ }
\Gamma(X,\mathcal{B})\xrightarrow{\ \beta\ }
\Gamma(X,\mathcal{C})
\]
\textit{is exact.}

As in the proof of \hyperref[fac:II:3:no:45:thm:2]{Theorem~2}, we may suppose that $X$ is the affine space $K^r$ and hence is irreducible. Set
\[
\mathcal{J}=\operatorname{Im}(\alpha)=\Ker(\beta).
\]
It remains to show that
\[
\alpha:\Gamma(X,\mathcal{A})\longrightarrow\Gamma(X,\mathcal{J})
\]
is surjective. By \hyperref[fac:II:3:no:45:thm:2:cor:1]{Corollary~1} there is a surjective homomorphism
\[
\varphi:\mathcal{O}_X^p\longrightarrow\mathcal{A},
\]
and \hyperref[fac:II:3:no:44:prop:7:cor:2]{Corollary~2 to Proposition~7} shows that
\[
\alpha\circ\varphi:
\Gamma(X,\mathcal{O}_X^p)\longrightarrow\Gamma(X,\mathcal{J})
\]
is surjective. It follows a fortiori that
\[
\alpha:\Gamma(X,\mathcal{A})\longrightarrow\Gamma(X,\mathcal{J})
\]
is surjective, as required.
